\documentclass[11pt,a4paper]{article}
\usepackage[margin=0.75in]{geometry}
\usepackage{enumitem}
\usepackage[hidelinks]{hyperref}
\setlist[itemize]{leftmargin=*,topsep=2pt,itemsep=1pt}
\pagestyle{empty}
\begin{document}

\begin{center}{\LARGE\bfseries Machine Learning Engineer – Model Training, Evaluation, and Deployment}\\[4pt]
Aniket Saxena | h20250169@pilani.bits-pilani.ac.in | github.com/aniketqw\end{center}

\section*{Summary}
Engineered end‑to‑end machine learning pipelines that ingest, clean, and validate data, and trained vision models to detect shots and extract keyframe features. Deployed models that output structured evidence in multiple formats and integrated LLM‑based semantic extraction into deterministic workflows.

\section*{Projects}
\begin{itemize}
  \item Implemented a shot‑detection pipeline using TransNet V2 or an adaptive detector to identify shot leaves and keyframes.
  \item Extracted CLIP or color features and sampled OCR from keyframes to enrich visual evidence.
  \item Built an adjacent chronological hierarchy that feeds an interactive timeline and node explorer, enabling evidence navigation.
  \item Configured the pipeline to output evidence in JSON, CSV, thumbnails, and a checksum ZIP, supporting downstream consumption.
  \item Designed a data processing pipeline that acquires, cleans, chunks, extracts, validates, and stores data in MongoDB, ensuring idempotent updates.
  \item Implemented semantic extraction using an LLM while keeping all other steps deterministic in Python.
  \item Integrated LangChain ReAct Agent with Ollama and MongoDB for agentic inference.
\end{itemize}

\section*{Skills}
\begin{itemize}
  \item TransNet V2
  \item CLIP
  \item OCR
  \item MongoDB
  \item Docker Compose
  \item LangChain
  \item Ollama
\end{itemize}

\end{document}